\documentclass[11pt,a4paper]{article}

\usepackage[utf8]{inputenc}
\usepackage[T1]{fontenc}
\usepackage[spanish,es-noquoting]{babel}
\usepackage[a4paper,margin=2.5cm]{geometry}
\usepackage{lmodern}
\usepackage{microtype}
\usepackage{xcolor}
\usepackage{booktabs}
\usepackage{longtable}
\usepackage{array}
\usepackage{enumitem}
\usepackage{graphicx}
\usepackage{float}
\usepackage{caption}
\usepackage{titlesec}
\usepackage{titletoc}
\usepackage{fancyhdr}
\usepackage{listings}
\usepackage{tikz}
\usetikzlibrary{arrows.meta, positioning}
\usepackage[hidelinks,colorlinks=true,linkcolor=azulref,urlcolor=azulref,citecolor=azulref]{hyperref}

\definecolor{azulref}{RGB}{20,80,150}
\definecolor{grisclaro}{RGB}{240,243,247}
\definecolor{grismedio}{RGB}{120,120,120}
\definecolor{verdeok}{RGB}{20,120,60}

\captionsetup{font=small,labelfont={bf,color=azulref}}

\lstdefinestyle{caja}{
  basicstyle=\ttfamily\footnotesize,
  backgroundcolor=\color{grisclaro},
  frame=single, rulecolor=\color{grismedio},
  breaklines=true, columns=fullflexible, keepspaces=true,
  showstringspaces=false, xleftmargin=4pt, xrightmargin=4pt,
  aboveskip=8pt, belowskip=8pt
}
\lstset{style=caja}

\titleformat{\section}{\Large\bfseries\color{azulref}}{\thesection.}{0.6em}{}
\titleformat{\subsection}{\large\bfseries}{\thesubsection}{0.6em}{}
\titlespacing*{\section}{0pt}{16pt}{8pt}
\titlecontents{section}
  [0em]
  {\bfseries\color{azulref}\addvspace{0.35em}}
  {\contentslabel{2.8em}}
  {\hspace*{-2.8em}}
  {\titlerule*[0.7pc]{.}\contentspage}
\titlecontents{subsection}
  [2.0em]
  {}
  {\contentslabel{3.2em}}
  {\hspace*{-3.2em}}
  {\titlerule*[0.7pc]{.}\contentspage}

\newcolumntype{L}[1]{>{\raggedright\arraybackslash}p{#1}}

\pagestyle{fancy}
\fancyhf{}
\rhead{\small ResearchFlow n8n}
\lhead{\small Proyecto final --- Daniel Cueto}
\rfoot{\small\thepage}
\lfoot{\small\color{grismedio} camba.tech}
\renewcommand{\headrulewidth}{0.4pt}
\renewcommand{\footrulewidth}{0.4pt}

\setlength{\parskip}{0.55em}
\setlength{\parindent}{0pt}
\newcommand{\code}[1]{\texttt{\small #1}}

\begin{document}

%============ PORTADA (formato oficial) ============
\begin{titlepage}
\centering
\vspace*{1.2cm}
{\large\bfseries PROYECTO FINAL}\\[0.4cm]
{\LARGE\bfseries AUTOMATIZACIÓN DE PROCESOS DE\\ NEGOCIO MEDIANTE n8n}\\[1.6cm]

\rule{0.82\textwidth}{1pt}\\[0.5cm]
{\Huge\bfseries\color{azulref} ResearchFlow}\\[0.35cm]
{\large Asistente de investigación automatizado y verificable con n8n}\\[0.3cm]
{\normalsize\color{verdeok}\bfseries Sistema desplegado y verificado en producción}\\[0.15cm]
{\small\color{grismedio}\url{https://n8n.camba.tech} $\cdot$ \url{https://chat.camba.tech}}\\[0.5cm]
\rule{0.82\textwidth}{1pt}\\[2.2cm]

\begin{tabular}{rl}
\textbf{Estudiante:} & Daniel Cueto Torrico\\[0.25cm]
\textbf{Programa:} & Automatización de Negocios con n8n\\[0.25cm]
\textbf{Institución:} & Embajada de los Estados Unidos\\[0.25cm]
\textbf{Docente:} & Ph.D. Milton Cayo Blanco\\[0.25cm]
\end{tabular}

\vspace{2cm}
{\normalsize 9 de julio de 2026}
\vfill
{\small\color{grismedio} n8n $\cdot$ PostgreSQL 16 $\cdot$ Gemini 2.5 Flash $\cdot$ Evolution API $\cdot$ QuickChart $\cdot$ Caddy $\cdot$ DigitalOcean}
\end{titlepage}

\tableofcontents
\newpage

%============ 1 ============
\section{Resumen Ejecutivo}
\textbf{Problema:} investigar un tema a fondo de forma manual (buscar fuentes, filtrar su calidad, extraer datos, construir gráficos y redactar un artículo) toma entre 4 y 6 horas por tema, y es propenso a sesgos, cifras sin respaldo y fuentes dudosas.

\textbf{Sector objetivo:} creador de contenido / analista independiente que produce artículos y análisis basados en investigación.

\textbf{Solución propuesta:} ResearchFlow, un asistente construido sobre n8n con \textbf{cuatro workflows}: (1) \emph{Investigación profunda} (webhook $\rightarrow$ Postgres $\rightarrow$ playbook editable $\rightarrow$ Gemini 2.5 Flash en 3 fases con \emph{Google Search grounding} $\rightarrow$ artículo HTML con gráficos por Gmail); (2) \emph{Ideas por WhatsApp} (AI Agent con memoria vía Evolution API); (3) \emph{Digest semanal}; y (4) \emph{Demo sin credenciales}.

\textbf{Herramientas:} n8n, PostgreSQL 16, Gemini 2.5 Flash + grounding, QuickChart, Gmail, Google Sheets, Evolution API (WhatsApp), Caddy (HTTPS) y Docker sobre un droplet de DigitalOcean con dominio propio \code{camba.tech}.

\textbf{Beneficios y resultados:} una investigación pasa de 4--6 horas a $\sim$10 minutos de ejecución más una revisión breve, con costo marginal de centavos. Por reglas duras del playbook nunca se inventan cifras ni fuentes. \textbf{El sistema fue desplegado y verificado end-to-end en producción}, con entrega real de artículos por correo (Sección~\ref{sec:deploy}).

%============ 2 ============
\section{Introducción}
La transformación digital permite convertir procesos repetitivos en flujos automatizados y medibles. El trabajo de investigación para contenido y análisis sigue siendo mayormente manual: se recibe una idea, se abren decenas de pestañas, se evalúan fuentes a ojo, se transcriben cifras a una hoja de cálculo y se redacta. Al agregar IA generativa sin proceso aparece un riesgo nuevo: respuestas fluidas que mezclan hechos comprobados con inferencias y cifras inventadas.

\textbf{Situación actual del negocio:} el analista invierte 4--6 h por tema, pierde ideas en canales dispersos y no tiene un archivo histórico reutilizable. \textbf{Necesidad de automatización:} industrializar la investigación con control de calidad y trazabilidad. \textbf{Importancia de la transformación digital:} pasar de un proceso artesanal a uno medible y escalable. \textbf{Justificación de n8n:} orquesta webhooks, IA, base de datos, APIs y notificaciones en flujos versionables y self-hosted, ideal para este proceso intelectual.

%============ 3 ============
\section{Requisitos Mínimos del Proyecto}
\renewcommand{\arraystretch}{1.3}
\begin{longtable}{L{5.3cm} L{9.7cm}}
\toprule
\textbf{Requisito obligatorio} & \textbf{Cómo se cumple en ResearchFlow}\\
\midrule \endhead
Un Trigger & Tres: Webhook \code{/webhook/researchflow}, Webhook \code{/webhook/researchflow-whatsapp} (Evolution) y Schedule Trigger (digest, lunes 8am).\\
Al menos tres integraciones externas & Seis: Gemini API (con grounding), Gmail, Google Sheets, Evolution API, QuickChart y PostgreSQL.\\
Base de datos o Google Sheets & Ambos: Postgres 16 (7 tablas) y Google Sheets (hoja \code{datasets}).\\
Proceso automatizado de comunicación & Artículo HTML por Gmail, avisos y conversación por WhatsApp, digest semanal.\\
Evidencias funcionales & Workflows desplegados y activos, ejecuciones verdes, artículo entregado por correo, registros en Postgres (Anexos).\\
Análisis económico & ROI y costos reales (Secciones~\ref{sec:roi} y~\ref{sec:econ}).\\
Documentación técnica & Informe, manual técnico, manual de usuario, esquema SQL y guía de despliegue.\\
\bottomrule
\end{longtable}

%============ 4 ============
\section{Planteamiento del Problema}
\textbf{Procesos manuales existentes:} búsqueda de fuentes, evaluación de calidad a ojo, extracción de cifras, armado de gráficos y redacción. \textbf{Tiempo invertido:} 4 a 6 horas por tema. \textbf{Errores frecuentes:} cifras sin fuente asociada, sesgo de selección de fuentes, mezcla de hechos con opiniones. \textbf{Costos asociados:} horas de trabajo del analista y riesgo reputacional por publicar datos indefendibles. \textbf{Impacto en la productividad:} solo 1--2 investigaciones por semana; las ideas llegan por canales dispersos y se pierden; no existe base histórica reutilizable. \textbf{Oportunidades de mejora:} automatizar captura, investigación metódica con IA, registro trazable y comunicación, con reglas de calidad estrictas que la IA debe cumplir.

%============ 5 ============
\section{Objetivos}
\subsection{Objetivo General}
Implementar con n8n un asistente de investigación automatizado que transforme un tema o pregunta en un artículo verificable con gráficos, reduciendo el tiempo de 4--6 horas a minutos, con trazabilidad completa de fuentes y comunicación automática por Gmail y WhatsApp.

\subsection{Objetivos Específicos}
\begin{itemize}[leftmargin=1.4em]
\item Analizar el proceso manual y codificar la metodología del usuario en un \emph{playbook} editable en base de datos.
\item Diseñar el flujo automatizado: 4 workflows en n8n.
\item Implementar la solución: Gemini en 3 fases con grounding, Postgres, Sheets, QuickChart, Gmail y WhatsApp.
\item Realizar pruebas unitarias e integrales.
\item Evaluar las mejoras obtenidas y desplegar en producción (droplet con HTTPS).
\end{itemize}

%============ 6 ============
\section{Marco Teórico}
\textbf{Automatización de procesos:} uso de tecnología para ejecutar tareas repetitivas con mínima intervención humana. ResearchFlow automatiza un proceso intelectual completo. \textbf{Transformación digital:} rediseño de procesos con herramientas digitales; la investigación pasa de artesanal a industrializada. \textbf{BPM:} el proceso ``investigar un tema'' se modeló en fases explícitas (playbook) y sus estados viven en la base de datos. \textbf{Integración de aplicaciones:} n8n conecta landing web, Postgres, Gemini, QuickChart, Gmail, Sheets y WhatsApp. \textbf{APIs REST:} el proyecto consume \code{generateContent} de Gemini, \code{sendText} de Evolution y QuickChart por URL. \textbf{Webhooks:} endpoints HTTP que reciben eventos y disparan workflows. \textbf{IA aplicada a negocios y grounding:} el grounding con Google Search ancla las respuestas de Gemini en resultados web reales y citables; combinado con reglas duras (no inventar cifras ni fuentes) la IA se vuelve auditable. \textbf{Plataforma n8n:} automatización visual self-hosted que conecta triggers, nodos de transformación, bases de datos, APIs y nodos de IA en flujos versionables en JSON.

%============ 7 ============
\section{Análisis del Negocio}
\subsection{Descripción del Negocio}
\textbf{Nombre/perfil:} analista de contenido independiente. \textbf{Sector económico:} creación de contenido, análisis e investigación aplicada. \textbf{Productos o servicios:} artículos con datos y gráficos, informes breves, backlog de temas. \textbf{Clientes objetivo:} el propio analista (asistente personal), extensible a equipos de contenido, consultores y estudiantes. \textbf{Problemas identificados:} tiempo alto por tema, riesgo por fuentes dudosas, pérdida de ideas y ausencia de archivo histórico.

\subsection{Proceso Seleccionado}
Investigación completa de un tema, desde la captura de la idea hasta la entrega del artículo verificable. Una persona envía un tema (y opcionalmente una pregunta) por la landing o WhatsApp. Sin pregunta: el sistema propone 3--5 preguntas investigables y las guarda como ideas. Con pregunta: responde ``en proceso'' y ejecuta la investigación en 3 fases con búsqueda web real, clasifica afirmaciones (hecho/hipótesis/opinión), extrae datasets con fuente, genera gráficos, guarda todo en Postgres y Sheets, y entrega el artículo por Gmail con aviso por WhatsApp.

%============ 8 ============
\section{Estudio de Viabilidad}
\subsection{Viabilidad Técnica}
Todos los nodos usados (Webhook, Set, Code, IF, Postgres, HTTP Request, Gmail, Google Sheets, Schedule Trigger, AI Agent con modelo Gemini y memoria) son estándar de n8n. El stack (Caddy, n8n, Postgres 16, Redis, Evolution API) corre con Docker Compose sobre un droplet de DigitalOcean. \textbf{La viabilidad técnica quedó confirmada: el sistema está desplegado y operativo en producción.}

\subsection{Viabilidad Económica}
Droplet cubierto por créditos de DigitalOcean; régimen posterior de 4~GB a \$24/mes. Gemini 2.5 Flash cuesta centavos por investigación. QuickChart, Vercel, Evolution API (open source) y las APIs de Google Workspace: \$0. Dominio \code{camba.tech}: $\sim$\$30/año. No hay licencias de software.

\subsection{Viabilidad Operativa}
El usuario solo llena un formulario o escribe por chat; el sistema hace el resto. La metodología es editable sin tocar los flujos (registro \code{playbook\_investigacion} en \code{app\_settings}). La supervisión humana se conserva: revisión del artículo antes de publicar y elementos marcados \code{requiere\_verificacion} explícitos.

%============ 9 ============
\section{Diseño de la Solución}
\subsection{Arquitectura General}

\begin{figure}[H]
\centering
\begin{tikzpicture}[font=\small,
  box/.style={rounded corners, draw=azulref, fill=grisclaro, minimum width=3.6cm, minimum height=1.0cm, align=center, thick, font=\footnotesize},
  big/.style={rounded corners, draw=azulref, fill=azulref, text=white, minimum width=8.2cm, minimum height=1.1cm, align=center, thick},
  arr/.style={-{Stealth[length=2.2mm]}, thick, draw=azulref}]
  \node[box] (landing) {Landing web\\(Vercel)};
  \node[box, right=1.4cm of landing] (wa) {WhatsApp / Simulador\\de chat};
  \node[big, below=1.2cm of landing, xshift=2.5cm] (n8n) {\textbf{n8n --- 4 workflows}\\[2pt]{\footnotesize droplet DigitalOcean $\cdot$ Docker $\cdot$ Caddy HTTPS $\cdot$ camba.tech}};
  \node[box, below=1.3cm of n8n, xshift=-3.9cm] (gem) {Gemini 2.5 Flash\\{\footnotesize + Search grounding}};
  \node[box, below=1.3cm of n8n] (pg) {PostgreSQL 16\\{\footnotesize 7 tablas research\_*}};
  \node[box, below=1.3cm of n8n, xshift=3.9cm] (out) {Gmail $\cdot$ Sheets\\{\footnotesize QuickChart $\cdot$ Evolution}};
  \draw[arr] (landing) -- (n8n);
  \draw[arr] (wa) -- (n8n);
  \draw[arr] (n8n) -- (gem);
  \draw[arr] (n8n) -- (pg);
  \draw[arr] (n8n) -- (out);
\end{tikzpicture}
\caption{Arquitectura de la solución}
\end{figure}

\subsection{Herramientas Utilizadas}
\renewcommand{\arraystretch}{1.2}
\begin{longtable}{L{4.6cm} L{10.4cm}}
\toprule
\textbf{Herramienta} & \textbf{Función}\\ \midrule \endhead
n8n (Docker) & Orquestación de los 4 workflows.\\
PostgreSQL 16 & Solicitudes, fuentes, evidencia, datasets, ideas y configuración (playbook).\\
Gemini 2.5 Flash + grounding & Investigación en 3 fases, generación de preguntas, refinado de ideas, digest.\\
QuickChart & Gráficos PNG por URL sin API key.\\
Gmail / Google Sheets (OAuth2) & Entrega del artículo y repositorio de datasets.\\
Evolution API v2.2 & Canal WhatsApp (captura de ideas, comandos y avisos).\\
Caddy + dominio \code{camba.tech} & Reverse proxy con HTTPS automático (Let's Encrypt).\\
Vercel & Landing con rewrite \code{/api/investigar} (evita CORS).\\
\bottomrule
\end{longtable}

%============ 10 ============
\section{Diseño del Flujo de Trabajo}

\begin{figure}[H]
\centering
\includegraphics[width=\textwidth]{output/pdf/img/fig_workflow_n8n.png}
\caption{Flujo de trabajo implementado en n8n (ejecución exitosa, todos los nodos en verde)}
\end{figure}

Workflow principal: \code{n8n\_workflow\_research\_production.json}. Recibe la solicitud por webhook, la registra en Postgres, carga el playbook y ejecuta la investigación en 3 fases (base y profundización con grounding, síntesis en JSON estricto), valida datasets, arma los gráficos y el artículo HTML, guarda la evidencia y entrega por Gmail, WhatsApp y Google Sheets.

\subsection{Descripción de Nodos}
\renewcommand{\arraystretch}{1.2}
\begin{longtable}{L{4.6cm} L{10.4cm}}
\toprule
\textbf{Nodo} & \textbf{Función}\\ \midrule \endhead
Webhook & Recepción de solicitudes (landing, WhatsApp o simulador).\\
Normalizar solicitud & Estandariza nombre, email, whatsapp, tema, pregunta y canal.\\
Postgres --- Registrar / Cargar playbook & INSERT con RETURNING id; lee la metodología editable de \code{app\_settings}.\\
IF --- Con pregunta & Bifurca: investigar o proponer preguntas.\\
Gemini --- Fase base / Profundizar / Síntesis & Investigación en 3 fases (grounding en 1--2; JSON estricto en 3).\\
Procesar resultado & Valida datasets (solo cifras con fuente), arma QuickChart, HTML y filas.\\
Gmail --- Enviar artículo & Envía el artículo HTML con gráficos incrustados.\\
Google Sheets --- Registrar datos & Una fila por punto de dato en la hoja \code{datasets}.\\
AI Agent --- Refinar idea & Agente con Gemini y memoria por número (flujo WhatsApp).\\
Evolution API & Notificaciones y conversación por WhatsApp.\\
Schedule Trigger & Digest semanal (lunes 8am).\\
\bottomrule
\end{longtable}

%============ 11 ============
\section{Implementación}
\label{sec:deploy}
\begin{enumerate}[leftmargin=1.6em]
\item \textbf{Configuración de credenciales:} Postgres, Gmail OAuth2 y Google Sheets OAuth2 (proyecto Google Cloud con redirect URI de n8n) y Google Gemini para el AI Agent. La configuración sensible (claves, IDs, datos del propietario) se centraliza en un \code{.env} que n8n inyecta a los workflows vía \code{\{\{ \$env.* \}\}}.
\item \textbf{Creación del flujo:} importación y activación de los 4 workflows.
\item \textbf{Integración de APIs:} Gemini (\code{generateContent} con \code{google\_search} en fases 1--2 y salida JSON estricta en la síntesis), QuickChart por URL y Evolution \code{sendText}.
\item \textbf{Webhooks:} la landing usa un rewrite de Vercel. Evolution envía el evento \path{messages.upsert} a \path{/webhook/researchflow-whatsapp}.
\item \textbf{Pruebas unitarias:} validación por nodo con \code{test\_data/} (normalización, INSERT, parseo del JSON de síntesis, URLs QuickChart).
\item \textbf{Pruebas integrales:} solicitud completa con entrega real de artículo por correo.
\item \textbf{Puesta en producción:} droplet Ubuntu 24.04, dominio \code{camba.tech} con HTTPS, workflows activos.
\end{enumerate}

\medskip
El sistema no quedó en diseño: fue desplegado y verificado en producción. Las siguientes subsecciones documentan el entorno real, la prueba end-to-end y los problemas de ingeniería resueltos.

\subsection{Entorno de producción}
\renewcommand{\arraystretch}{1.2}
\begin{longtable}{L{4.4cm} L{10.6cm}}
\toprule \textbf{Componente} & \textbf{Valor real}\\ \midrule \endhead
Servidor & Droplet DigitalOcean, Ubuntu 24.04, IP \code{165.227.203.200}.\\
Dominio & \code{camba.tech} (subdominios \code{n8n.}, \code{evo.}, \code{chat.}).\\
n8n & \url{https://n8n.camba.tech} --- \code{HTTP/2 200} con HTTPS válido.\\
Contenedores & caddy, postgres (healthy), redis, n8n, evolution, simulador --- todos \emph{running}.\\
Base de datos & 7 tablas creadas; playbook sembrado en \code{app\_settings}.\\
\bottomrule
\end{longtable}

\subsection{Prueba integral end-to-end (exitosa)}
Se lanzó una investigación real por el webhook de producción. Respuesta inmediata:
\begin{lstlisting}
{"success":true,"modo":"investigacion","estado":"en_proceso",
 "request_id":2,"mensaje":"Investigacion iniciada. El articulo llegara
 a tu correo en unos minutos."}
\end{lstlisting}
En $\sim$1m~35s corrieron las 3 fases de Gemini, se guardó la evidencia en Postgres, se registraron los datos en Google Sheets y \textbf{el artículo HTML con su gráfico llegó al correo}. En la base de datos, las solicitudes quedan con estado \code{completado\_con\_busqueda}, lo que confirma que el pipeline completo (incluida la búsqueda web y el envío por Gmail) se ejecutó.

\begin{figure}[H]
\centering
\includegraphics[width=\textwidth]{output/pdf/img/fig_ejecucion_verde.png}
\caption{Ejecución exitosa del workflow en n8n (panel \emph{Executions}, todos los nodos en verde)}
\end{figure}

\subsection{Problemas resueltos durante el despliegue}
\renewcommand{\arraystretch}{1.25}
\begin{longtable}{L{4.6cm} L{10.4cm}}
\toprule \textbf{Problema} & \textbf{Solución aplicada}\\ \midrule \endhead
Imagen de Evolution inaccesible & Se cambió al namespace oficial vigente \code{evoapicloud/evolution-api:v2.2.3}.\\
Placeholders y claves inline en los JSON & Refactor a variables de entorno (\code{\{\{ \$env.* \}\}}) leídas del \code{.env}: sin secretos en archivos versionados.\\
Llamada interna frágil a \code{localhost} & El flujo de WhatsApp usa la URL pública del webhook (\code{RESEARCH\_WEBHOOK\_URL}).\\
Modelo de IA no disponible (503) & Se fijó el modelo estable \code{gemini-2.5-flash} en el AI Agent y los nodos HTTP.\\
Restricción de WhatsApp (Baileys) & Ante la inestabilidad del canal WhatsApp no oficial, se construyó un \textbf{simulador de chat} que ejercita el mismo flujo real.\\
\bottomrule
\end{longtable}

\subsection{Canal de chat: simulador desplegado}
El canal conversacional se expone además como un \textbf{simulador de chat} desplegado en \url{https://chat.camba.tech}, que reproduce el bot (mismo prompt y modelo Gemini) y cuyo comando \code{investigar} dispara la investigación profunda real. Es una entrada estable e independiente del WhatsApp no oficial.

\begin{figure}[H]
\centering
\includegraphics[width=0.55\textwidth]{output/pdf/img/fig_chat_simulador.png}
\caption{Simulador de chat: el asistente refina una idea vaga en un tema con preguntas investigables}
\end{figure}

%============ 12 ============
\section{Resultados}
\begin{table}[H]
\centering\renewcommand{\arraystretch}{1.3}
\begin{tabular}{L{5cm} L{4cm} L{4.5cm}}
\toprule
\textbf{Indicador} & \textbf{Antes (manual)} & \textbf{Después (ResearchFlow)}\\ \midrule
Tiempo por investigación & 4 a 6 horas & $\sim$10 min + revisión breve\\
Investigaciones por semana & 1 a 2 & 5 a 10\\
Cifras sin fuente & Frecuentes & 0 (sin fuente no se grafica)\\
Errores / trazabilidad & Dispersa & 100\% en Postgres, clasificada\\
Costo operativo por tema & Horas de trabajo & Centavos de API\\
Productividad & Baja & 4--5$\times$ mayor\\
\bottomrule
\end{tabular}
\caption{Comparación de resultados antes y después de la automatización}
\end{table}
La mejora central no es solo velocidad: cada artículo queda \textbf{auditable} (qué fuente respalda qué hecho) y el conocimiento se acumula en una base consultable.

%============ 13 ============
\section{Retorno de Inversión (ROI)}
\label{sec:roi}
\begin{itemize}[leftmargin=1.4em]
\item Investigaciones/mes: 16. Tiempo manual: 5~h c/u; con ResearchFlow: $\sim$0.67~h. \textbf{Ahorro: 4.33~h por investigación.}
\item \textbf{Tiempo ahorrado:} $\sim$4.3~h/día, $\sim$21~h/semana, $\sim$69~h/mes.
\item \textbf{Costo proceso manual:} 16$\times$5$\times$\$6 = \$480/mes. \textbf{Costo automatizado:} \$25/mes (droplet 4~GB + Gemini).
\item Valor hora del analista: \$6. Inversión inicial: $\sim$\$240 (una vez).
\end{itemize}
\begin{lstlisting}
Ahorro mensual          = 69.3 h x $6            = $415.80
Costo mensual sistema   =                          $25.00
Beneficio neto mensual  = 415.80 - 25.00         = $390.80
ROI mensual             = 390.80 / 25 x 100      = 1,563%
Recuperacion inversion  = 240 / 390.80           = ~3 semanas
\end{lstlisting}
Incluso recortando los supuestos a la mitad, el beneficio neto supera con holgura el costo del sistema.

%============ 14 ============
\section{Análisis Económico}
\label{sec:econ}
\subsection{Costos}
\renewcommand{\arraystretch}{1.2}
\begin{longtable}{L{10.5cm} L{4cm}}
\toprule \textbf{Ítem} & \textbf{Costo/mes}\\ \midrule \endhead
Infraestructura --- Droplet (con créditos DigitalOcean) & \$48 $\rightarrow$ \$0 efectivo\\
Hosting --- Droplet 4~GB (post-créditos) & \$24\\
APIs --- Gemini 2.5 Flash (uso personal) & $<$ \$1\\
Licencias (n8n, Postgres, Evolution: open source) & \$0\\
Servicios externos (QuickChart, Vercel, Gmail/Sheets) & \$0\\
Dominio \code{camba.tech} & $\sim$\$2.5\\
\bottomrule
\end{longtable}

\subsection{Beneficios}
\begin{itemize}[leftmargin=1.4em]
\item \textbf{Reducción de errores:} cero cifras sin fuente publicadas.
\item \textbf{Ahorro de tiempo:} $\sim$69 horas al mes.
\item \textbf{Incremento de productividad:} 4--5$\times$ más investigaciones con el mismo tiempo.
\item \textbf{Escalabilidad:} metodología editable en datos; el estilo se ajusta sin tocar los flujos.
\item \textbf{Mejor experiencia:} entrega automática por Gmail, WhatsApp/chat y digest semanal.
\end{itemize}

%============ 15 ============
\section{Conclusiones}
\begin{enumerate}[leftmargin=1.6em]
\item \textbf{Automatización:} n8n permite automatizar un proceso intelectual completo (investigación en 3 fases con grounding, clasificación de evidencia y entrega multicanal) sin intervención humana.
\item \textbf{Productividad y transformación digital:} la producción crece 4--5$\times$ y la investigación pasa de artesanal a un servicio real en internet (Docker, HTTPS, dominio propio, webhooks públicos).
\item \textbf{Impacto económico:} con \$25/mes de costo total, el ROI mensual supera el 1{,}500\% bajo supuestos conservadores; la inversión se recupera en $\sim$3 semanas.
\item La IA generativa aporta valor cuando se somete a reglas verificables: el playbook con reglas duras convierte a Gemini de un riesgo de alucinaciones en un asistente auditable.
\end{enumerate}

%============ 16 ============
\section{Recomendaciones}
\begin{itemize}[leftmargin=1.4em]
\item Integrar fuentes de datos primarias por API (Our World in Data, bancos centrales, INE) como herramientas del flujo.
\item Implementar RAG sobre las investigaciones pasadas (tablas \code{research\_*}).
\item Migrar el canal WhatsApp a la API oficial de Meta (Cloud API) para estabilidad de producción.
\item Incorporar un segundo modelo como verificador cruzado de los hechos \code{requiere\_verificacion}.
\item Programar respaldos automáticos de Postgres (\code{pg\_dump}) además del snapshot del droplet.
\end{itemize}

%============ 17 ============
\section{Bibliografía}
\begin{enumerate}[leftmargin=2.2em,label={[\arabic*]}]
\item n8n GmbH, \emph{n8n Documentation}. \url{https://docs.n8n.io/}
\item Google, \emph{Gemini API --- Grounding with Google Search}.\par \url{https://ai.google.dev/gemini-api/docs/google-search}
\item Google, \emph{Gemini Developer API --- Pricing}. \url{https://ai.google.dev/gemini-api/docs/pricing}
\item QuickChart, \emph{Documentation}. \url{https://quickchart.io/documentation/}
\item Evolution API, \emph{Documentation}. \url{https://doc.evolution-api.com/}
\item Docker Inc., \emph{Docker Compose Documentation}. \url{https://docs.docker.com/compose/}
\item DigitalOcean, \emph{Droplets Documentation}.\par \url{https://docs.digitalocean.com/products/droplets/}
\item PostgreSQL Global Development Group, \emph{PostgreSQL 16 Documentation}. \url{https://www.postgresql.org/docs/16/}
\item Caddy, \emph{Automatic HTTPS}. \url{https://caddyserver.com/docs/automatic-https}
\end{enumerate}

%============ 18 ============
\section{Rúbrica de Evaluación}
\renewcommand{\arraystretch}{1.25}
\begin{longtable}{L{7.5cm} r L{5cm}}
\toprule \textbf{Criterio} & \textbf{Pts} & \textbf{Evidencia}\\ \midrule \endhead
Identificación del problema empresarial & 10 & Sec.~4 y 7\\
Justificación de la automatización & 10 & Sec.~8 y 13\\
Diseño de la solución & 15 & Sec.~9 y 10, Figura 1--2\\
Implementación funcional en n8n & 20 & Sec.~11, Anexos\\
Integraciones externas & 15 & Gemini, Gmail, Sheets, Evolution, QuickChart, Postgres\\
Innovación y creatividad & 10 & 3 fases + playbook editable + anti-alucinación\\
Análisis económico y ROI & 10 & Sec.~13 y 14\\
Calidad del informe técnico & 5 & Este informe + manuales\\
Presentación y defensa & 5 & Demo en vivo (chat + correo real)\\
\midrule \textbf{Total} & \textbf{100} & \\
\bottomrule
\end{longtable}

%============ 19 ============
\section{Anexos}
Evidencias funcionales del sistema en producción:

\begin{figure}[H]
\centering
\includegraphics[width=0.92\textwidth]{output/pdf/img/fig_evolution.png}
\caption{Anexo A --- Integración de WhatsApp: instancia \code{researchflow} conectada en Evolution API}
\end{figure}

\begin{figure}[H]
\centering
\includegraphics[width=0.92\textwidth]{output/pdf/img/fig_dns.png}
\caption{Anexo B --- Despliegue en dominio propio: registros DNS \code{n8n}, \code{evo} y \code{chat} apuntando al droplet}
\end{figure}

La guía operativa está en \path{DOCUMENTACION.md}. El paquete incluye además los workflows JSON, el esquema SQL, la infraestructura en \path{deploy/} y el simulador de chat.

\vspace{0.8cm}
\begin{center}\small\color{grismedio}\rule{0.5\textwidth}{0.4pt}\\[0.25cm]
ResearchFlow --- desplegado en \url{https://n8n.camba.tech} $\cdot$ Daniel Cueto Torrico $\cdot$ 9 de julio de 2026
\end{center}

\end{document}
